% LaTeX template for Artifact Evaluation V20201122
%
% Prepared by 
% * Grigori Fursin (cTuning foundation, France) 2014-2020
% * Bruce Childers (University of Pittsburgh, USA) 2014
%
% See examples of this Artifact Appendix in
%  * SC'17 paper: https://dl.acm.org/citation.cfm?id=3126948
%  * CGO'17 paper: https://www.cl.cam.ac.uk/~sa614/papers/Software-Prefetching-CGO2017.pdf
%  * ACM ReQuEST-ASPLOS'18 paper: https://dl.acm.org/citation.cfm?doid=3229762.3229763
%
% (C)opyright 2014-2020
%
% CC BY 4.0 license
%

\documentclass{sigplanconf}

\usepackage{hyperref}

\begin{document}

\special{papersize=8.5in,11in}

%%%%%%%%%%%%%%%%%%%%%%%%%%%%%%%%%%%%%%%%%%%%%%%%%%%%
% When adding this appendix to your paper, 
% please remove above part
%%%%%%%%%%%%%%%%%%%%%%%%%%%%%%%%%%%%%%%%%%%%%%%%%%%%

\appendix
\section{Artifact Appendix}

%%%%%%%%%%%%%%%%%%%%%%%%%%%%%%%%%%%%%%%%%%%%%%%%%%%%%%%%%%%%%%%%%%%%%
\subsection{Abstract}

This artifact supports two experiments for the paper
``Personalized Federated Hyperdimensional Computing via Shared--Personal
Prototype Memory Decomposition'': (1) controlled-system reproduction of
Tables I, II, and III, and (2) a six-dataset benchmark comparing HoRU,
HyperFeel, and FedHDC. The table experiment generates its controlled fixture
locally and produces the latency and payload tables. The benchmark experiment
acquires public inputs, builds local caches, runs the shared protocol, and
writes validation reports.

\subsection{Artifact check-list (meta-information)}

{\em Obligatory. Use just a few informal keywords in all fields applicable to your artifacts
and remove the rest. This information is needed to find appropriate reviewers and gradually 
unify artifact meta information in Digital Libraries.}

{\small
\begin{itemize}
  \item {\bf Algorithm: } personalized federated HDC, low-rank shared/personal prototype decomposition
  \item {\bf Program: } Python package and command-line workflows
  \item {\bf Compilation: } none; interpreted Python
  \item {\bf Binary: } optional CUDA-enabled PyTorch
  \item {\bf Model: } prototype memories and low-rank bases built at run time
  \item {\bf Data set: } controlled-system fixture, UCI-HAR, ISOLET, FEMNIST, WISDM, LEAF Synthetic, NinaPro DB1
  \item {\bf Run-time environment: } Python 3.11+, Linux
  \item {\bf Hardware: } CPU for Tables I--III; one CUDA GPU recommended for the full benchmark
  \item {\bf Execution: } command line
  \item {\bf Metrics: } bootstrap latency, recurring-round latency, upload payload, round-25 accuracy
  \item {\bf Output: } CSV, JSON, Markdown reports
  \item {\bf Experiments: } table reproduction; six-dataset benchmark and validation
  \item {\bf How much disk space required (approximately)?: } table workflow under 1 GB; benchmark inputs and outputs require tens of GB
  \item {\bf How much time is needed to prepare workflow (approximately)?: } a few minutes for the Python environment; dataset acquisition depends on downloads
  \item {\bf How much time is needed to complete experiments (approximately)?: } tables in minutes; quick benchmark in tens of minutes; full benchmark is GPU-dependent and longer
  \item {\bf Publicly available?: } yes
  \item {\bf Code licenses (if publicly available)?: } repository terms
  \item {\bf Data licenses (if publicly available)?: } upstream dataset licenses
  \item {\bf Workflow framework used?: } none
  \item {\bf Archived (provide DOI)?: } none assigned in this repository snapshot
\end{itemize}
}

%%%%%%%%%%%%%%%%%%%%%%%%%%%%%%%%%%%%%%%%%%%%%%%%%%%%%%%%%%%%%%%%%%%%%
\subsection{Description}

\subsubsection{How to access}

The repository root contains the artifact. The active code paths are:
\begin{itemize}
  \item \texttt{src/horu\_artifact/experiments/tables123.py} for Tables I--III;
  \item \texttt{artifact/scripts/run\_cuda\_reconstruction\_suite.py} for the
        six-dataset benchmark;
  \item \texttt{reference\_results/cuda\_suite/} for immutable
        reference reports;
  \item \texttt{configs/accuracy\_full.yaml},
        \texttt{configs/accuracy\_quick.yaml}, and
        \texttt{configs/datasets.yaml} for benchmark configuration.
\end{itemize}

\subsubsection{Hardware dependencies}

The table workflow is CPU-only. The six-dataset reference reports were
generated with CUDA; reproducing the full benchmark is practical with one
CUDA-capable GPU. The code still runs on CPU, but the full benchmark may be
slow.

\subsubsection{Software dependencies}

Python 3.11 or newer is required. The package dependencies are
\texttt{torch >= 2.2}, \texttt{PyYAML >= 6.0}, and \texttt{scipy >= 1.11}.
Tests use \texttt{pytest >= 8}.

\subsubsection{Data sets}

Tables I--III use a generated controlled-system fixture and do not require an
external dataset. The benchmark uses six public datasets with pinned
input-acquisition scripts: UCI-HAR, ISOLET, FEMNIST, WISDM, Synthetic, and
NinaPro DB1. Short dataset notes are under \texttt{docs/datasets/}.

\subsubsection{Models}

No pretrained neural network is required. The artifact constructs prototype
memories, shared bases, and personal bases during execution.

%%%%%%%%%%%%%%%%%%%%%%%%%%%%%%%%%%%%%%%%%%%%%%%%%%%%%%%%%%%%%%%%%%%%%
\subsection{Installation}

Create a Python environment and install the package:

\begin{verbatim}
python3 -m venv .venv
source .venv/bin/activate
pip install -e '.[test]'
\end{verbatim}

For CUDA benchmark runs, install the appropriate CUDA-enabled PyTorch wheel
before the editable install.

%%%%%%%%%%%%%%%%%%%%%%%%%%%%%%%%%%%%%%%%%%%%%%%%%%%%%%%%%%%%%%%%%%%%%
\subsection{Experiment workflow}

\paragraph{Tables I, II, and III.}
Generate the controlled-system fixture, then run the table workflow:

\begin{verbatim}
horu-artifact prepare-data controlled-systems --data-root data
horu-artifact reproduce-tables \
  --data-root data \
  --output results/table_reproduction \
  --warmup 5 \
  --repeats 30 \
  --threads 1
\end{verbatim}

\paragraph{Six-dataset benchmark.}
First acquire the public inputs:

\begin{verbatim}
python3 artifact/scripts/acquire_uci_har_prototype.py \
  --source-root /data/horu/uci_har \
  --archive /data/downloads/UCI_HAR_Dataset.zip

python3 artifact/scripts/acquire_isolet_prototype.py \
  --source-root /data/horu/isolet \
  --download-dir /data/downloads/isolet

python3 artifact/scripts/prepare_femnist_reconstruction.py \
  --source-root /data/horu/femnist

python3 artifact/scripts/acquire_wisdm_reconstruction.py \
  --source-root /data/horu/wisdm \
  --outer-archive /data/downloads/WISDM.zip

python3 artifact/scripts/prepare_synthetic_reconstruction.py \
  --source-root /data/horu/synthetic

python3 artifact/scripts/acquire_ninapro_db1_reconstruction.py \
  --source-root /data/horu/ninapro \
  --download-dir /data/downloads/ninapro
\end{verbatim}

Then run the benchmark:

\begin{verbatim}
python3 artifact/scripts/run_cuda_reconstruction_suite.py \
  --uci-har-source-root /data/horu/uci_har \
  --isolet-raw-source-root /data/horu/isolet \
  --femnist-source-root /data/horu/femnist \
  --wisdm-source-root /data/horu/wisdm \
  --synthetic-source-root /data/horu/synthetic \
  --ninapro-db1-source-root /data/horu/ninapro \
  --output-dir results_reconstruction/cuda_suite
\end{verbatim}

This command builds local caches, runs the shared protocol, and writes the
validation report.

%%%%%%%%%%%%%%%%%%%%%%%%%%%%%%%%%%%%%%%%%%%%%%%%%%%%%%%%%%%%%%%%%%%%%
\subsection{Evaluation and expected results}

\paragraph{Tables I, II, and III.}
The table workflow must generate:
\begin{itemize}
  \item \texttt{table1.csv}
  \item \texttt{table2.csv}
  \item \texttt{table3.csv}
  \item \texttt{raw\_timings.csv}
  \item \texttt{environment.json}
  \item \texttt{result.json}
\end{itemize}
The expected qualitative outcome is that HoRU shows lower recurring-round
latency and much smaller upload payload than FedHDC and HyperFeel.

\paragraph{Six-dataset benchmark.}
The immutable reference report is under
\texttt{reference\_results/cuda\_suite/}. It can be checked with:

\begin{verbatim}
python3 artifact/scripts/verify_reconstruction_suite.py \
  --manifest artifact/manifests/reconstruction_cuda_suite_v1.json \
  --suite-output reference_results/cuda_suite
\end{verbatim}

The expected round-25 accuracies in the committed reference reports are:
\begin{itemize}
  \item UCI-HAR: HoRU 97.66, HyperFeel 98.08, FedHDC 95.25
  \item ISOLET: HoRU 88.14, HyperFeel 90.59, FedHDC 93.85
  \item FEMNIST: HoRU 68.54, HyperFeel 67.78, FedHDC 57.26
  \item WISDM: HoRU 56.69, HyperFeel 51.11, FedHDC 6.95
  \item Synthetic: HoRU 78.00, HyperFeel 73.19, FedHDC 71.09
  \item NinaPro DB1: HoRU 75.68, HyperFeel 61.48, FedHDC 30.86
  \item Unweighted mean: HoRU 77.45, HyperFeel 73.71, FedHDC 59.21
\end{itemize}

%%%%%%%%%%%%%%%%%%%%%%%%%%%%%%%%%%%%%%%%%%%%%%%%%%%%%%%%%%%%%%%%%%%%%
\subsection{Experiment customization}

The main benchmark configuration files are:
\begin{itemize}
  \item \texttt{configs/accuracy\_quick.yaml}
  \item \texttt{configs/accuracy\_full.yaml}
  \item \texttt{configs/datasets.yaml}
\end{itemize}
Reviewers can switch between quick and full benchmark settings, change output
directories, or run only the table workflow if GPU resources are unavailable.

%%%%%%%%%%%%%%%%%%%%%%%%%%%%%%%%%%%%%%%%%%%%%%%%%%%%
% When adding this appendix to your paper, 
% please remove below part
%%%%%%%%%%%%%%%%%%%%%%%%%%%%%%%%%%%%%%%%%%%%%%%%%%%%

\end{document}
